\documentclass[25pt, a0paper, portrait, colspace=8mm]{tikzposter}
\usepackage[beamer]{prettytex/base}
\usepackage{prettytex/math}
% \usepackage[backend=biber,style=numeric]{biblatex}
% \usepackage[toc, acronym, style=long3col, indexonlyfirst=true, nogroupskip=true]{glossaries}
\usepackage[backend=biber, citestyle=numeric, bibstyle=authoryear, hyperref=true, sorting=none, maxbibnames=99]{biblatex}
\usepackage{csquotes}
\usepackage[nameinlink]{cleveref}
\usepackage{dsfont}

\usepackage{times}
\usepackage{xfrac}
\usepackage{bm}

\ifdefined\norm
  \renewcommand{\norm}[1]{\left\lVert#1\right\rVert_{\scriptscriptstyle 2}}
\else
  \newcommand{\norm}[1]{\left\lVert#1\right\rVert_{\scriptscriptstyle 2}}
\fi

\newcommand{\midrule}{\hline}

\addbibresource{../../literature/sources.bib}

\title{Real-time interaction with an electrophysiological cancer cell model}
\author{
  Peter Waldert\textsuperscript{1,2},
  Theresa Rienmüller\textsuperscript{1},
  Christian Baumgartner\textsuperscript{1} and
  Sonja Langthaler\textsuperscript{1}
}
\date{\today}
\institute{
  \textsuperscript{1}Institute of Health Care Engineering with European Testing Center of Medical Devices, TU Graz \\
  \textsuperscript{2}Institute of Computer Graphics and Knowledge Visualisation, TU Graz
}

\definecolor{TUGraz}{HTML}{f70146}
% \definecolor{TUGraz}{HTML}{e4154b}
\definecolor{Accent}{RGB}{136, 58, 234}
\definecolor{AccentGreen}{RGB}{0, 128, 0}
\definecolor{AccentBlue}{RGB}{0, 0, 255}

\definecolorstyle{TUGrazColorStyle}{
  \colorlet{colorTwo}{TUGraz}
}{
  % Background Colors
  \colorlet{backgroundcolor}{white}
  \colorlet{framecolor}{TUGraz}
  % Title Colors
  \colorlet{titlefgcolor}{black}
  \colorlet{titlebgcolor}{TUGraz}
  % Block Colors
  \colorlet{blocktitlebgcolor}{Accent}
  \colorlet{blocktitlefgcolor}{black}
  \colorlet{blockbodybgcolor}{colorTwo}
  \colorlet{blockbodyfgcolor}{black}
  % Innerblock Colors
  \colorlet{innerblocktitlebgcolor}{white}
  \colorlet{innerblocktitlefgcolor}{black}
  \colorlet{innerblockbodybgcolor}{white}
  \colorlet{innerblockbodyfgcolor}{black}
  % Note colors
  \colorlet{notefgcolor}{black}
  \colorlet{notebgcolor}{white}
  \colorlet{notefrcolor}{white}
}

\defineblockstyle{TUGrazBlockStyle}{
  titlewidthscale=1, bodywidthscale=1,
  titleoffsetx=0pt, titleoffsety=0pt, bodyoffsetx=0pt, bodyoffsety=0pt,
  bodyverticalshift=0pt, roundedcorners=0, linewidth=0.2cm,
  titleinnersep=6mm, bodyinnersep=6mm, titleleft
}{
  \begin{scope}[line width=\blocklinewidth, rounded corners=\blockroundedcorners]
    \ifBlockHasTitle %
      \draw[draw=none]
          (blockbody.south west) rectangle (blocktitle.north east);
      \draw[color=blocktitlebgcolor, solid, thin]
          (blocktitle.south west) -- (blocktitle.south east);%
    \else
      \draw[draw=none]
          (blockbody.south west) rectangle (blockbody.north east);
    \fi
  \end{scope}
}

\usetheme{Simple}
\usetitlestyle{Empty}
\usecolorstyle{TUGrazColorStyle}
\useblockstyle{TUGrazBlockStyle}

\settitle{
  \vspace{0.4cm}
  \hspace{0.5cm}
  \vbox{
    \centering
    \color{titlefgcolor}
    {\bfseries \Huge \sc \@title \par}
    \vspace*{0.8em}
    {\LARGE \@author \par}
    \vspace*{0.6em}
    {\Large \@institute}  
    \vspace*{-1.8cm}
  }
}

\begin{document}
  \hypersetup{urlcolor=Accent, linkcolor=AccentBlue, citecolor=AccentBlue}
  \renewcommand{\operatorcolor}{AccentBlue}

  \maketitle[width=\linewidth]

  \node[anchor=west] at ([xshift=1.8cm, yshift=-2.5cm]TP@title.west) {\includegraphics[height=4.0cm]{../logos/tugraz.pdf}};
  % \node[anchor=east] at ([xshift=-14cm, yshift=-6cm]TP@title.east) {\includegraphics[height=2.5cm]{../logos/hce.pdf}};
  \node[anchor=east] at ([xshift=2.8cm, yshift=-2.5cm]TP@title.east) {\includegraphics[height=5.4cm]{../logos/biotechmed-graz.eps}};
  % \node[] at (33.1in/2-10cm, -46.8in/2+4cm) {\includegraphics[width=15cm]{../logos/biotechmed-graz.eps}};

  \begin{columns}
    \column{0.25}
    \block{\sc Introduction}{
      % We extend the previous work \cite{2021-A549-model} by a faster, in-browser implementation in Rust \cite{2014-rust} using compilation to WebAssembly.

      We improve on the A549 electrophysiological cancer cell model introduced in \cite{2021-A549-model,2024-calcium-channels}, combining numerical methods with an efficient implementation to reduce simulation time to a level where it is feasible for live interaction.

      \includegraphics[width=\linewidth]{../../figures/cell-by-langthaler-et-al.png}
      \textit{Schematic depiction of the cell, source: \cite{2024-calcium-channels}}
      \vspace*{2mm}

      We accelerated the simulation with adaptive timestepping and a highly efficient implementation in the Rust programming language.
      The corresponding inverse problem is expressed as a quadratic program, solving it within milliseconds.
      We introduce a visualisation approach of the entire model: a live simulation dashboard available online, running directly in the browser.
      % The entire source code is freely available on GitHub and reusable through three different channels: the simulation interface (powered by compilation to WebAssembly), the Rust linkable library implementation and a Python package.
      \vspace*{-2.4cm}
    }

    \column{0.75}
    \block{\sc Live, In-Browser Cell Simulation Interface\vspace*{-0.115em}}{
      \includegraphics[width=\linewidth, trim={0 2.5cm 0 0}, clip]{../../figures/interface-a4.pdf}
      \textit{This interface is available online via \href{https://in-silico.hce.tugraz.at/}{in-silico.hce.tugraz.at}, try it out!}
      \vspace*{-1.5cm}
    }
  \end{columns}

  \begin{columns}
    \column{0.5}
    \block{\sc Model}{
    The whole cell current $I: T \to \R$ over time $t \in T \subset \R^+$ is the sum of all individual channel contributions $I_k, k \in \{1, ..., M\}$ over $M \in \N$ channel types
    \begin{equation*}
      I(t) := \sum_{k=1}^{M} N_k I_k(t) = \sum_{k=1}^{M} N_k g_k p_{o,k} \left(V(t)-E_k\right)\,.
      \label{eq:current}
    \end{equation*}
    % where $N_k$ is the number of channels of type $k \in \{1, ..., M\}$, $g_k$ is the respective ion channel's conductivity, $p_{o, k} \in [0, 1]$ is the probability of observing the channel in a state where an ion current can flow (``open states''), $V: T \to \R$ is the voltage across the membrane and $E_k \in \R$ the reversal potential.

    At each time step, the next state $\vec{s}_{k,n+1} \in [0, 1]^{N_{s,k}}$ of the $k$-th channel type is obtained by
    \begin{equation*}
      \vec{s}_{k,n+1} = H_{k}\left(V(t_n), \vec{C}(t_n), t_n\right) \vec{s}_{k,n}\,,
      \quad\text{with}\quad
      t_n := \sum_{i=0}^n (\Delta t)_i\,.
    \end{equation*}
    % where $\vec{s}_{k,n} \in [0, 1]^{N_{s,k}}$ is the state vector of ion channel type $k$ at the $n$-th time step, $H_{k}\left(V, \vec{C}, t_n\right) \in [0, 1]^{N_{s,k} \times N_{s,k}}$ the transition matrix for type $k$ with $\sum_{j=1}^{N_{s,k}} \{H_k\}_{i,j} = 1 \;\forall\,i$, $V(t_n)$ the voltage across the membrane at time $t_n$ and $\vec{C}(t_n) \in \R^4$ the concentrations of Kalium, Calcium, Sodium and Chlorine at time $t_n$.
    \vspace*{-2.8cm}
    }
    \block{\sc Adaptive Timestepping}{
      In order to accelerate the simulation in areas where there is little change to the dynamics, we choose an adaptive step size $(\Delta t)_{n}$ based on
      \begin{equation*}
        (\Delta t)_{n+1} = (\Delta t)_{n} \left(\frac{\Delta^{\rm tol}}{\sum_{k=1}^{M} N_k \norm{\vec{s}_{k,n+1} - \vec{s}_{k,n}}}\right)^{1/2}\,.
        \label{eq:dt-heuristic}
      \end{equation*}

      % \begin{figure}
      In order to find the optimal $\Delta^{\rm tol} \in \R^+$ (the allowed state change in between steps):
      \vspace{4mm}

      \includegraphics[width=\linewidth]{../../figures/results/delta-tolerance.pdf}

      \textit{Relative change of the average timestep $\Delta t$ (in \col{AccentBlue}{blue}), simulation runtime (in \col{Accent}{violet}) and step acceptance rate (in \col{AccentGreen}{green}) when varying the delta tolerance $\Delta^{\rm tol}$ on a log-scale. All three quantities were normalized from their individual extent to $[0, 1]$. The most effective $\Delta^{\rm tol}$ is arguably on the order of $10^{-4}$.}
      % \label{figure:delta-tolerance}
      % \end{figure}
      \vspace*{-1.0cm}
    }
    % \block{\sc Pointwise Error between Simulation and Measurements}{
    %   \includegraphics[width=\linewidth]{../../figures/results/simulation-error.pdf}
    %   % \caption{Pointwise error between the measured current $\vec{I}_{\rm meas}$ and the simulated current $\vec{I}$, showing potential systematic problems within the computational model's representation of the ion channels as compared to the experimental results.}
    %   \vspace*{-3.0cm}
    % }

    \column{0.5}
    \block{\sc Formulation as a Quadratic Program\vspace*{-0.155em}}{
      After simulation of the individual channel currents, we want to find
      \begin{equation*}
        \vec{N}_{\rm opt} = \argmin_{\vec{N} \in \N_0^M} \frac{1}{2} \norm{R \vec{N} - \vec{I}_{\rm meas}}^2\,.
        \label{eq:minimization}
      \end{equation*}
      the number of channels per type, with $\vec{I}_{\rm meas} \in \R^{N_t}$ the experimentally measured current and $R \in \R^{N_t \times M}$ the matrix of all currents $I_k$ per channel type.
      Letting $\vec{d} := \vec{I}_{\rm meas}$ for brevity,
      $$\vec{N}_{\rm opt} \approx \argmin_{\vec{x} \in \R_+^M} f(\vec{x}) = \argmin_{\vec{x} \in \R_+^M} \sfrac{1}{2} \norm{R \vec{x} - \vec{d}}^2\,,$$
      with cost function $f: \R^M \to \R^+$, which we manipulate to
      \begin{align*}
        f(\vec{x}) & = \sfrac{1}{2} (R \vec{x} - \vec{d})^T (R \vec{x} - \vec{d})                                                          \\
                   & = \sfrac{1}{2} \left(\vec{x}^T R^T R \vec{x} - \vec{x}^T R^T \vec{d} - \vec{d}^T R \vec{x} + \vec{d}^T \vec{d}\right) \\
                   & = \sfrac{1}{2} \left(\vec{x}^T P \vec{x} + \vec{x}^T \vec{q} + \vec{q}^T \vec{x}\right) + \mathcal{O}(1)              \\
                   & = \sfrac{1}{2} \vec{x}^T P \vec{x} + \vec{q}^T \vec{x} + \mathcal{O}(1)
      \end{align*}
      where we let $P := R^T R \in \R^{M \times M}$ and $\vec{q} := -R^T \vec{d} \in \R^M$ and leave out the constant $\vec{d}^T \vec{d}$ as $\mathcal{O}(1)$.
      We can express the nonnegativity constraint $\vec{x} \ge \vec{0}$ as an equality constraint using a slack variable $\vec{s} \in \R_+^M$,
      $$-\vec{x} + \vec{s} = \vec{0} \quad\Leftrightarrow\quad A \vec{x} + \vec{s} = \vec{b}\,\quad \text{with}\quad A = -\mathds{1} \in \R^{M \times M} \;\;\text{and}\;\; \vec{b} = \vec{0} \in \R^M\,.$$
      This leaves us with a constrained \textit{quadratic program},
      \begin{align*}
        \min_{\vec{x} \in \R^M}\; & \sfrac{1}{2}\, \vec{x}^T P \vec{x} + \vec{q}^T \vec{x},  \\
        s.t.\;                    & A \vec{x} + \vec{s} = \vec{b}\,,\; \vec{s} \in \R_+^M\,.
      \end{align*}

      % We solve the quadratic problem in this exact form using Clarabel \cite{2024-clarabel}.
      % Note that in Clarabel notation, the slack variable is to be taken as an element of the nonnegativity cone.

      The integer solution can then be obtained from rounding,
      $$\vec{N}_{\rm opt} = \lfloor \vec{x} \rceil \in \N_0^M\,.$$
      \vspace*{-3.0cm}
    }
    \block{\sc Comparison of Optimization Methods}{
      % \begin{table*}[ht]
      % \centering
      The different approaches were evaluated on the G0 cell cycle phase with the activation voltage protocol. Runtime estimates were obtained on an Intel\texttrademark i7-5600U CPU.
      \vspace{.5cm}

      \begin{tabular}{lcrr}
        \hline
        \textbf{Algorithm}                          & \textbf{Abbreviation} & \textbf{Runtime} / ms & \textbf{RMSE} / pA \\
        \midrule
        Particle Swarm Optimization                 & PSO                   & 22571                 & 27.69              \\
        Gradient Descent + More Thuente             & GD                    & 18924                 & 32.34              \\
        Limited-Memory BFGS + Hager Zhang           & LBFGS                 & 4845                  & 32.20              \\
        Non-Negative Least Squares \cite{1997-nnls} & NNLS                  & 318                   & 28.00              \\
        Quadratic Program                           & QP                    & 18                    & 28.13              \\
      \end{tabular}
      % \label{table:optimization-comparison}
      % \end{table*}
      \vspace*{-3.0cm}
    }
  \end{columns}

  \block{\sc Conclusion}{
    Efforts to further enhance and complete the first electrophysiological cancer cell model simulation were fruitful, resulting in a new library implementation, an improved inverse problem solution technique and a live in-browser simulation dashboard.
    Numerically, the stability of the simulation varies greatly with the time step state change tolerance $\Delta^{\rm tol}$, this could be improved using a higher-order integration scheme.
    There are also many ideas for further customizations that would improve the live simulation dashboard, opening up future usage perspectives.
    \vspace*{-1.0cm}
  }

  \block{\sc References}{
    \vspace*{-0.2cm}
    \begin{minipage}{\linewidth}
      \begin{center}\mbox{}
        \printbibliography[heading=none]
      \end{center}
    \end{minipage}
  }
\end{document}
